% =============================================================================
% Chapter 14: Robust AI — ML Systems Lecture Slides (Volume II)
% =============================================================================
% Covers: silent failure modes, distribution shift, adversarial attacks,
% data poisoning, defense strategies, and the robustness-accuracy trade-off.
%
% IMAGES NEEDED (SVGs converted to PDF):
%   - silent-failure-modes.pdf    - robustness-framework.pdf
%   - distribution-shift-types.pdf - adversarial-attack-fgsm.pdf
%   - robustness-tax.pdf          - sdc-at-scale.pdf
%   - data-poisoning.pdf          - defense-layers.pdf
% =============================================================================
\documentclass[aspectratio=169, 12pt]{beamer}
\usepackage{../../assets/beamerthememlsys}

\mlsyssetup{
  volume       = {Volume II},
  chapter      = {Chapter 14},
  logo         = {../../assets/img/logo-mlsysbook.png},
  instlogo     = {../../assets/img/logo-harvard.png},
  chaptertitle = {Robust AI},
}

% --- Fonts ---
\usepackage{fontspec}
\setsansfont{Helvetica Neue}[
  BoldFont={Helvetica Neue Bold},
  ItalicFont={Helvetica Neue Italic},
  BoldItalicFont={Helvetica Neue Bold Italic},
]
\setmonofont{JetBrains Mono}[Scale=0.85]

% --- Packages ---
\usepackage{booktabs}
\usepackage{amsmath}

% --- Image paths ---
\graphicspath{
  {images/}
}

% --- Chapter-specific macros ---
\newcommand{\eps}{\varepsilon}

% --- Helper: safe image include ---
\newcommand{\safeimg}[2][width=\textwidth,keepaspectratio]{%
  \IfFileExists{images/#2}{\includegraphics[#1]{#2}}{%
    \IfFileExists{#2}{\includegraphics[#1]{#2}}{%
      \fbox{\parbox[c][2.5cm][c]{0.85\linewidth}{\centering\footnotesize\textcolor{midgray}{[Missing image]}}}%
    }%
  }%
}

% --- Section count for navigation ---
\setcounter{mlsystotalsections}{8}

\title{Robust AI}
\author{Vijay Janapa Reddi}
\institute{Harvard University}
\date{}

\begin{document}

% =============================================================================
% TITLE SLIDE
% =============================================================================
\mlsystitle{Robust AI}{Making the Invisible Visible}{cover_robust_ai.png}

% =============================================================================
% LEARNING OBJECTIVES
% =============================================================================
\begin{frame}{Learning Objectives}
\note{
% -- NARRATE: What to SAY while showing this slide
[2 min] Walk through each objective. Emphasize that robustness is a systems
property, not just an algorithmic feature. ``This chapter makes the invisible
visible --- silent failures that no alert catches.''

% -- ENGAGE: Specific question, prediction, or task for THIS slide
Ask: ``How would you know if your deployed model was silently failing right now?''
Cold-call one student. [Expected: most will admit they would not know.]

% -- FLEX: [CORE] or [OPTIONAL] + contingency
[CORE] This slide sets the chapter scope.
IF SHORT: Read objectives 1, 2, and 5 aloud; let students scan the rest.
}

\footnotesize
\begin{enumerate}\setlength\itemsep{0pt}
  \item Classify robustness challenges into \textbf{environmental shifts}, \textbf{adversarial attacks}, and \textbf{data poisoning}
  \item Explain why ML systems \textbf{fail silently} and quantify the \textbf{robustness--accuracy trade-off}
  \item Evaluate adversarial attack techniques (\textbf{FGSM}, \textbf{PGD}, \textbf{C\&W}) and their threat models
  \item Apply statistical methods (\textbf{MMD}, \textbf{PSI}, \textbf{KS tests}) to detect distribution shifts
  \item Construct \textbf{multi-layered defenses} combining adversarial training, certified defenses, and input sanitization
  \item Integrate robustness with \textbf{compute and accuracy trade-offs} across the ML lifecycle
\end{enumerate}

\end{frame}

\begin{frame}{Visual Language}
\note{
% -- NARRATE: What to SAY while showing this slide
[1 min] Point to each card. ``In this chapter, red dominates --- robustness
failures are errors and bottlenecks. Green marks the healthy paths we are
trying to protect.''

% -- FLEX: [OPTIONAL] + contingency
[OPTIONAL] Skip if students have seen this in a previous lecture.
}

\small
Throughout this course, colors carry meaning:

\vspace{0.3cm}
\begin{columns}[T]
  \begin{column}{0.45\textwidth}
    \begin{mlsyscard}{computestroke}
    \textbf{Blue} --- Compute / Processing\\
    {\footnotesize GPU ops, forward/backward pass, inference}
    \end{mlsyscard}
    \vspace{0.15cm}
    \begin{mlsyscard}{datastroke}
    \textbf{Green} --- Data / Memory\\
    {\footnotesize Data flow, caches, healthy paths}
    \end{mlsyscard}
  \end{column}
  \begin{column}{0.45\textwidth}
    \begin{mlsyscard}{routingstroke}
    \textbf{Orange} --- Routing / Scheduling\\
    {\footnotesize Load balancers, batch windows}
    \end{mlsyscard}
    \vspace{0.15cm}
    \begin{mlsyscard}{errorstroke}
    \textbf{Red} --- Error / Cost / Bottleneck\\
    {\footnotesize Loss, decode phase, waste}
    \end{mlsyscard}
  \end{column}
\end{columns}
\end{frame}



% =============================================================================
\section{Silent Failures}
% =============================================================================

\begin{frame}{The Silent Failure Problem}
\note{
% -- LINK: What prior concept connects to this slide
Ch 13 covered intentional attacks. This chapter opens with a different threat:
the model fails on its own, without any adversary, because the world changed.

% -- NARRATE: What to SAY while showing this slide
[3 min] Point to the diagram: ``Traditional software crashes loudly with stack
traces. ML models produce confident but wrong outputs --- no exception, no alert,
no crash.'' The autonomous vehicle scenario makes this visceral: ``The model
kept driving at full speed because it was 98\% confident the white trailer was sky.''

% -- ENGAGE: Specific question, prediction, or task for THIS slide
Ask: ``What is more dangerous --- a system that crashes or one that keeps running
while silently wrong?'' Give 10 seconds. [Expected: silently wrong, because no
one knows to intervene.]

% -- WARN: What students will get wrong on THIS topic
Common error: students think monitoring accuracy on a test set catches silent
failures. Correct framing: the test set is i.i.d. --- it cannot detect failures
on out-of-distribution inputs the model encounters in production.

% -- FLEX: [CORE] or [OPTIONAL] + contingency
[CORE] This is the chapter thesis --- do not skip.
IF SHORT: State ``ML fails silently'' and show one line of the diagram.
}

% --- Layout: FULL-WIDTH diagram ---
\centering
\safeimg[width=0.92\textwidth,height=4.8cm,keepaspectratio]{silent-failure-modes.pdf}

\vspace{0.15cm}
\mlsysalert{Key:}{No code changed. No crash. No alert. The \emph{world} changed.}

\end{frame}

\begin{frame}{Defining Robust AI}
\note{
% -- LINK: What prior concept connects to this slide
The previous slide showed that ML fails silently. This slide defines what
``robust'' means formally --- a measurable systems property, not a vague goal.

% -- NARRATE: What to SAY while showing this slide
[2 min] Read the definition card aloud. Then point to the table: ``Standard
accuracy measures average-case on i.i.d.\ test data. Robustness measures
worst-case under perturbation.'' Anchor on the key number: ``Epsilon = 8/255
is invisible to humans but drops ResNet-50 from 76\% to 11\%.''

% -- WARN: What students will get wrong on THIS topic
Common error: students think robustness = higher accuracy. Correct framing:
95\% test accuracy can coexist with 100\% vulnerability to targeted
perturbations. They measure fundamentally different things.

% -- FLEX: [CORE] or [OPTIONAL] + contingency
[CORE] The formal definition structures all subsequent slides.
IF SHORT: Read the definition and the epsilon=8/255 number only.
}

\footnotesize
\begin{mlsyscard}{crimson}
\textbf{Robust AI:} The measurable systems property that predictions remain valid ---
within specified error bounds --- under distribution shift, adversarial perturbation,
and hardware/software faults.
\end{mlsyscard}

\vspace{0.1cm}
\renewcommand{\arraystretch}{1.1}
{\scriptsize
\begin{tabular}{@{}lll@{}}
  \toprule
  & \textbf{Standard Accuracy} & \textbf{Robustness} \\
  \midrule
  Measures     & Average-case (i.i.d.\ test) & Worst-case (adversarial / OOD) \\
  Perturbation & None                         & $\eps = 8/255$ ($L_\infty$) \\
  Accuracy     & 76\% (ResNet-50)             & 11\% after single bit flip \\
  Failure mode & Visible in metrics           & Silent, high-confidence \\
  \bottomrule
\end{tabular}
}

\vspace{0.1cm}
\mlsysconcept{Distinction:}{95\% test accuracy can coexist with 100\% vulnerability to targeted perturbations.}

\end{frame}

% =============================================================================
\section{Real-World Failures}
% =============================================================================

\begin{frame}{Cloud: AWS S3 Outage (2017)}
\note{
% -- LINK: What prior concept connects to this slide
We defined robustness formally. Now we ground it in real failures --- starting
with infrastructure failure that cascades into ML service outage.

% -- NARRATE: What to SAY while showing this slide
[2 min] ``A single mistyped command during S3 maintenance cascaded through
US-East-1, disabling 150+ services for 4 hours. Alexa --- 40 million devices ---
went completely dark.'' Point to the card: ``ML services assume infrastructure
availability as an invariant, not a probabilistic guarantee.''

% -- WARN: What students will get wrong on THIS topic
Common error: students treat infrastructure failures as ``not an ML problem.''
Correct framing: ML services inherit all infrastructure fragility AND add
model-specific failure modes on top.

% -- FLEX: [OPTIONAL] or [CORE] + contingency
[OPTIONAL] This is context-setting for the Tesla case study.
IF SHORT: State the \$150M number and Alexa impact in one sentence, move on.
}

\footnotesize
\begin{columns}[T]
  \begin{column}{0.55\textwidth}
    \textbf{What happened:}
    \begin{itemize}\setlength\itemsep{0pt}
      \item Engineer mistyped a maintenance command
      \item US-East-1 region: 150+ services down for \textbf{4 hours}
      \item Estimated \textbf{\$150M} in losses
    \end{itemize}

    \vspace{0.1cm}
    \textbf{ML Impact:}
    \begin{itemize}\setlength\itemsep{0pt}
      \item Alexa (40M+ devices): completely unresponsive
      \item Voice recognition: 200--500 ms $\to$ total failure
    \end{itemize}
  \end{column}
  \begin{column}{0.42\textwidth}
    \begin{mlsyscard}{errorstroke}
    \textbf{Systems Lesson}\\[0.05cm]
    {\scriptsize ML services assume infrastructure availability as an \emph{invariant}, not a probabilistic guarantee. A single human error disables the entire ML fleet.}
    \end{mlsyscard}
  \end{column}
\end{columns}

\end{frame}

\begin{frame}{Edge: Tesla Autopilot (2016)}
\note{
% -- LINK: What prior concept connects to this slide
S3 was infrastructure failure. Tesla is a model-level robustness failure ---
the system worked correctly on inputs it had seen, but failed fatally on a
novel scenario.

% -- NARRATE: What to SAY while showing this slide
[3 min] Walk through the bullet list: ``Camera: white trailer against bright
sky --- missed. Radar: trailer classified as overhead sign --- ignored. Fusion
logic: discarded both valid signals. Result: no braking at 74 mph.'' Key
technical insight: ``Both sensors received valid data. The decision boundary
failed on an input absent from training. This is not a sensor failure --- it
is a robustness failure.''

% -- ENGAGE: Specific question, prediction, or task for THIS slide
Ask: ``What would have prevented this crash at the systems level?'' Give 10
seconds. [Expected: redundant decision logic, uncertainty quantification that
triggers fallback when confidence is low on novel inputs.]

% -- WARN: What students will get wrong on THIS topic
Common error: students blame the camera or radar hardware. Correct framing:
both sensors worked correctly --- the fusion logic's decision boundary had a
gap on out-of-distribution inputs.

% -- FLEX: [CORE] or [OPTIONAL] + contingency
[CORE] This is the anchor case study for edge robustness failures.
IF SHORT: Focus on the ``both sensors received valid data'' insight.
}

\footnotesize
\begin{columns}[T]
  \begin{column}{0.55\textwidth}
    \textbf{Fatal crash:} Tesla Model S at 74 mph
    \begin{itemize}\setlength\itemsep{0pt}
      \item Camera: white trailer vs.\ bright sky $\to$ missed
      \item Radar: trailer classified as overhead sign $\to$ ignored
      \item Fusion logic: discarded both valid signals
      \item Result: no braking, underride collision
    \end{itemize}

    \vspace{0.05cm}
    \begin{mlsyscard}{errorstroke}
    {\scriptsize Both sensors received valid data. The \textbf{decision boundary} failed on an input absent from training.}
    \end{mlsyscard}
  \end{column}
  \begin{column}{0.42\textwidth}
    \begin{mlsyscard}{routingstroke}
    \textbf{Edge Fragility}\\[0.05cm]
    {\scriptsize Edge devices lack fallback redundancy. No secondary cluster to route to. Must degrade gracefully within power and memory envelope.}
    \end{mlsyscard}
  \end{column}
\end{columns}
\mlsyscite{NTSB, 2017}

\end{frame}

\begin{frame}{Silent Data Corruption at Scale}
\note{
% -- LINK: What prior concept connects to this slide
Tesla showed an edge device failing on novel input. This slide shows a different
failure mode: hardware silently corrupting model weights at datacenter scale.

% -- NARRATE: What to SAY while showing this slide
[3 min] Walk through the probability math in the diagram: ``P(at least 1 SDC)
= 1 - (1-p) to the N. At 10K GPUs with Meta's reported per-device rate, you
get 63\% probability of at least one silent data corruption per hour.'' Key
punchline: ``A single bit flip can drop ResNet-50 from 76\% to 11\%.''

% -- ENGAGE: Specific question, prediction, or task for THIS slide
Ask: ``Does ECC memory solve this?'' Give 5 seconds.
[Expected: yes. Correct answer: no --- ECC catches correctable errors but not
all silent data corruptions, especially in logic paths.]

% -- WARN: What students will get wrong on THIS topic
Common error: students think ECC memory eliminates silent corruption. Correct
framing: ECC catches DRAM bit flips but not logic-path SDC, GPU register
corruption, or interconnect errors.

% -- FLEX: [CORE] or [OPTIONAL] + contingency
[CORE] The 63\% per hour number is the key quantitative anchor for scale risk.
IF SHORT: State the probability and the 76\%-to-11\% impact number.
}

% --- Layout: FULL-WIDTH diagram ---
\centering
\safeimg[width=0.92\textwidth,height=4.8cm,keepaspectratio]{sdc-at-scale.pdf}

\vspace{0.1cm}
\mlsysalert{Scale transforms negligible per-device risk into near-certain system-level failure.}{}

\end{frame}

% --- ACTIVE LEARNING 1: Predict ---
\begin{frame}{Predict: What Makes ML Failures Unique?}
\note{
% -- LINK: What prior concept connects to this slide
Three real-world failures (S3, Tesla, SDC) showed different robustness threats.
This prediction exercise asks students to synthesize those into general causes.

% -- NARRATE: What to SAY while showing this slide
[2 min] Read the scenario aloud: ``95\% accuracy drops to 72\% over 3 months.
No code changed. No alerts fired.'' Give 60 seconds of silent writing. Do NOT
reveal the answer yet. Ask 2--3 students to share.

% -- ENGAGE: Specific question, prediction, or task for THIS slide
Students must name 3 possible causes. [Expected answers: distribution shift,
data pipeline corruption, hardware SDC, upstream feature changes, label drift.]

% -- FLEX: [CORE] or [OPTIONAL] + contingency
[CORE] This active learning moment primes the unified framework on the next slide.
IF SHORT: Reduce writing time to 30 seconds; still cold-call 2 students.
}

\centering
\vspace{0.8cm}
{\Large\bfseries Think--Write--Share}

\vspace{0.5cm}
{\large A deployed model's accuracy drops from 95\% to 72\%\\
over 3 months. No code changed. No alerts fired.\\[0.3cm]
\alert{Name 3 possible causes.}}

\vspace{0.5cm}
{\normalsize Write them down. \textcolor{midgray}{(60 seconds)}}

\vspace{0.3cm}
{\small\textcolor{lightgray}{Then compare with a neighbor.}}

\end{frame}

% =============================================================================
\section{Unified Framework}
% =============================================================================

\begin{frame}{A Unified Robustness Framework}
\note{
% -- LINK: What prior concept connects to this slide
Students just predicted causes of silent degradation. This framework organizes
all those causes (and their defenses) into a unified structure.

% -- NARRATE: What to SAY while showing this slide
[3 min] Walk through the diagram: ``Four threat sources on the left ---
environmental shift, adversarial attack, data poisoning, hardware fault. Five
defense layers on the right --- input sanitization, robust training, drift
monitoring, data validation, uncertainty quantification. The model sits in the
center.'' Key punchline: ``No single defense addresses all four threats.''
Connect to prior chapters: ``Security (Ch 13) locks the front door. Fault
tolerance (Ch 7) keeps training alive. Robustness faces what both miss.''

% -- WARN: What students will get wrong on THIS topic
Common error: students think adversarial training alone makes a model robust.
Correct framing: adversarial training addresses one threat (perturbations) but
not distribution shift, hardware faults, or data poisoning.

% -- FLEX: [CORE] or [OPTIONAL] + contingency
[CORE] This framework is the chapter's organizing structure.
IF SHORT: Name the four threats and five defenses without elaborating each.
}

% --- Layout: FULL-WIDTH diagram ---
\centering
\safeimg[width=0.92\textwidth,height=4.8cm,keepaspectratio]{robustness-framework.pdf}

\vspace{0.1cm}
\mlsysinsight{Pattern:}{Every threat demands a distinct defense. Multi-layered strategy is non-negotiable.}

\end{frame}

\mlsysfocus{Robust AI}{%
$\text{Robustness} \neq \text{Accuracy}$\\[0.3cm]
A model with 95\% test accuracy can be\\
100\% vulnerable to targeted perturbations%
}

\begin{frame}{From Previous Chapters}
\note{
% -- LINK: What prior concept connects to this slide
The framework showed four threat sources. This slide maps prior chapter knowledge
to what robustness adds on top.

% -- NARRATE: What to SAY while showing this slide
[2 min] Point to the left column: ``You already know fault tolerance, security,
serving, and operations.'' Point to the right: ``Robustness adds graceful degradation,
adversarial resilience, and data integrity.'' Read the crimson card: ``A system that
is secure but fragile is useless. Robustness closes the loop.''

% -- FLEX: [OPTIONAL] + contingency
[OPTIONAL] Bridging slide. IF SHORT: Skip --- the framework slide covered this.
}

\scriptsize
\begin{columns}[T]
  \begin{column}{0.48\textwidth}
    \textbf{What students already know:}
    \begin{itemize}\setlength\itemsep{0pt}
      \item \textcolor{computestroke}{Fault Tolerance}: Checkpointing, recovery
      \item \textcolor{datastroke}{Security}: Access control, threat models
      \item \textcolor{routingstroke}{Serving}: Batching, pipeline parallelism
      \item \textcolor{errorstroke}{Ops}: Monitoring, SLAs
    \end{itemize}

    \vspace{0.05cm}
    \textbf{What robustness adds:}
    \begin{itemize}\setlength\itemsep{0pt}
      \item \alert{Graceful degradation} under shift
      \item \alert{Adversarial resilience} at inference
      \item \alert{Data integrity} in training pipelines
    \end{itemize}
  \end{column}
  \begin{column}{0.48\textwidth}
    \begin{mlsyscard}{crimson}
    \textbf{The Robustness Gap}\\[0.05cm]
    {\scriptsize A system that is secure but fragile is useless. One that tolerates faults but not drift will silently degrade. Robustness closes the loop.}
    \end{mlsyscard}

    \vspace{0.1cm}
    \begin{mlsyscard}{errorstroke}
    {\scriptsize \textbf{Scale amplifies risk:} 175B-param model on 8+ GPU nodes $\to$ 8$\times$ fault surface. INT8 quantization reduces robustness margin.}
    \end{mlsyscard}
  \end{column}
\end{columns}

\end{frame}

% =============================================================================
\section{Distribution Shift}
% =============================================================================

\begin{frame}{Three Types of Distribution Shift}
\note{
% -- LINK: What prior concept connects to this slide
The framework listed environmental shift as threat source 1. This slide decomposes
it into three types with different detection strategies.

% -- NARRATE: What to SAY while showing this slide
[3 min] Walk through the diagram: ``Covariate shift: P(X) changes. Label shift:
P(Y) changes. Concept drift: P(Y|X) changes.'' Key distinction: ``Covariate and
label shift are detectable from input statistics. Concept drift requires
ground-truth comparison --- much harder in production.''

% -- ENGAGE: Specific question, prediction, or task for THIS slide
Ask: ``Which type is hardest to detect and why?'' Give 10 seconds.
[Expected: concept drift --- input monitoring sees nothing unusual.]

% -- WARN: What students will get wrong on THIS topic
Common error: students assume all drift is detectable from input statistics.
Correct framing: concept drift is invisible to feature monitoring because P(X)
is unchanged.

% -- FLEX: [CORE] or [OPTIONAL] + contingency
[CORE] The three-type taxonomy structures detection methods on the next slide.
}

% --- Layout: FULL-WIDTH diagram ---
\centering
\safeimg[width=0.92\textwidth,height=4.8cm,keepaspectratio]{distribution-shift-types.pdf}

\vspace{0.1cm}
\mlsysconcept{Critical distinction:}{Concept drift changes P(Y$|$X) --- detectable \emph{only} with ground-truth feedback.}

\end{frame}

\begin{frame}{Detecting Drift: The Statistical Arsenal}
\note{
% -- LINK: What prior concept connects to this slide
Three drift types need detection tools. This slide provides the statistical
monitoring toolkit.

% -- NARRATE: What to SAY while showing this slide
[3 min] Walk through the table: ``KS test is fast but univariate. PSI bins
features and compares proportions. MMD works in high dimensions but is
compute-heavy.'' Point to the green card: ``PSI detects shifts 3--6 weeks
before accuracy falls below SLA.'' Point to the red card: ``Concept drift is
the blind spot --- invisible to all input-side methods.''

% -- WARN: What students will get wrong on THIS topic
Common error: students deploy one statistical test and assume full coverage.
Correct framing: KS covers univariate, MMD covers joint distribution, but
neither catches concept drift.

% -- FLEX: [CORE] or [OPTIONAL] + contingency
[CORE] The 3--6 week early warning is the key operational insight.
IF SHORT: Point to the table and green card only.
}

\small
\renewcommand{\arraystretch}{1.15}
{\footnotesize
\begin{tabular}{@{}llll@{}}
  \toprule
  \textbf{Method} & \textbf{What It Detects} & \textbf{Strength} & \textbf{Limitation} \\
  \midrule
  \textbf{KS Test}  & Univariate shift   & Fast ($O(n \log n)$)  & 1D only \\
  \textbf{PSI}      & Feature stability   & Easy to interpret     & Binning artifacts \\
  \textbf{MMD}      & Joint distribution  & High-dimensional      & Compute-heavy \\
  \textbf{Chi-sq}   & Categorical shift   & Standard test         & Discrete features \\
  \bottomrule
\end{tabular}
}

\vspace{0.2cm}
\begin{columns}[T]
  \begin{column}{0.48\textwidth}
    \begin{mlsyscard}{datastroke}
    \textbf{Early Warning}\\[0.05cm]
    {\footnotesize PSI detects shifts \textbf{3--6 weeks} before accuracy falls below SLA thresholds, enabling proactive retraining.}
    \end{mlsyscard}
  \end{column}
  \begin{column}{0.48\textwidth}
    \begin{mlsyscard}{errorstroke}
    \textbf{Blind Spot}\\[0.05cm]
    {\footnotesize Concept drift ($P(Y|X)$ changes) is \emph{invisible} to input monitoring. Requires ground-truth comparison.}
    \end{mlsyscard}
  \end{column}
\end{columns}

\end{frame}

\begin{frame}{Worked Example: Is the World Changing?}
\note{
% -- LINK: What prior concept connects to this slide
The statistical arsenal showed tools. This worked example applies them to a
concrete production scenario.

% -- NARRATE: What to SAY while showing this slide
[2 min] Walk through the table row by row: ``Baseline mean 0.50, observed 0.55.
Difference = 0.05, standard error = 0.009, Z-score = 5.5 sigma, p-value less
than 0.001.'' Key insight: ``A 5\% shift sounds small, but with 1,000 samples
it is 5.5 standard errors --- statistically certain.''

% -- WARN: What students will get wrong on THIS topic
Common error: students confuse statistical significance with practical significance.
Correct framing: 5.5 sigma confirms the shift is real, but the engineering question
is whether 0.05 shift matters for model accuracy.

% -- FLEX: [OPTIONAL] or [CORE] + contingency
[OPTIONAL] Reinforces the detection methods. IF SHORT: State the 5.5-sigma result.
}

\small
\textbf{Scenario:} Baseline feature mean = 0.50. Last 1,000 requests: mean = 0.55.

\vspace{0.15cm}
\renewcommand{\arraystretch}{1.15}
{\footnotesize
\begin{tabular}{@{}llr@{}}
  \toprule
  \textbf{Step} & \textbf{Calculation} & \textbf{Value} \\
  \midrule
  Difference   & $0.55 - 0.50$               & 0.05 \\
  Std.\ Error  & $0.3 / \sqrt{1000}$         & 0.009 \\
  Z-score      & $0.05 / 0.009$              & \textbf{5.5$\sigma$} \\
  P-value      & ---                          & \textbf{$< 0.001$} \\
  \bottomrule
\end{tabular}
}

\vspace{0.15cm}
\begin{mlsyscard}{errorstroke}
\textbf{Verdict: Confirmed Regression.} A shift of 0.05 with 1,000 samples has $< 0.1\%$ probability of being noise. Trigger automated response: fallback model or emergency retraining.
\end{mlsyscard}

\end{frame}

% =============================================================================
\section{Adversarial Attacks}
% =============================================================================

\begin{frame}{Adversarial Attacks: Exploiting Decision Boundaries}
\note{
% -- LINK: What prior concept connects to this slide
Distribution shift is accidental. Adversarial attacks are intentional --- they
weaponize the same gradient used for training.

% -- NARRATE: What to SAY while showing this slide
[3 min] Point to the FGSM equation in the diagram: ``The same gradient used for
training is weaponized for attack. Training minimizes loss; attack maximizes it.
One gradient step is enough.'' Anchor on the number: ``Epsilon = 8/255 is
invisible to humans but devastating to models.''

% -- WARN: What students will get wrong on THIS topic
Common error: students think adversarial vulnerability is a bug to fix. Correct
framing: it is an inherent property of high-dimensional decision boundaries ---
the geometry of the space enables it.

% -- FLEX: [CORE] or [OPTIONAL] + contingency
[CORE] FGSM is the foundation for all adversarial attack discussion.
IF AHEAD: ``Why does one gradient step suffice when training needs thousands?''
}

% --- Layout: FULL-WIDTH diagram ---
\centering
\safeimg[width=0.92\textwidth,height=4.8cm,keepaspectratio]{adversarial-attack-fgsm.pdf}

\vspace{0.1cm}
\mlsysalert{Not a bug:}{Adversarial vulnerability is an inherent property of high-dimensional decision boundaries.}

\end{frame}

\begin{frame}{Attack Taxonomy}
\note{
% -- LINK: What prior concept connects to this slide
FGSM showed the mechanism. This slide organizes all adversarial attacks into
four categories of increasing realism.

% -- NARRATE: What to SAY while showing this slide
[2 min] Walk through the table: ``Gradient-based requires white-box access ---
unlikely in production. Transfer-based is the realistic threat: craft on a
surrogate, transfer to production with 30--70\% success. Physical patches are
the most visceral: stickers on stop signs.''

% -- WARN: What students will get wrong on THIS topic
Common error: students focus on white-box attacks because the math is cleaner.
Correct framing: transfer and physical attacks are the production threats ---
they require zero access to the target model.

% -- FLEX: [OPTIONAL] or [CORE] + contingency
[OPTIONAL] Extends the FGSM slide with taxonomy context.
IF SHORT: Focus on transfer and physical rows only.
}

\scriptsize
\renewcommand{\arraystretch}{1.0}
\begin{tabular}{@{}llll@{}}
  \toprule
  \textbf{Category} & \textbf{Mechanism} & \textbf{Access} & \textbf{Key Example} \\
  \midrule
  \textbf{Gradient}    & $\nabla_x J(\theta,x,y)$ & White-box    & FGSM, PGD, JSMA \\
  \textbf{Optimization} & Min.\ perturbation    & White-box    & C\&W (broke distillation) \\
  \textbf{Transfer}    & Surrogate $\to$ target    & \textcolor{errorstroke}{\textbf{Black-box}} & 30--70\% success rate \\
  \textbf{Physical}    & Real-world patches        & \textcolor{errorstroke}{\textbf{No access}} & Stop sign stickers \\
  \bottomrule
\end{tabular}

\vspace{0.1cm}
\begin{columns}[T]
  \begin{column}{0.48\textwidth}
    \begin{mlsyscard}{errorstroke}
    {\scriptsize \textbf{Transfer Attacks}: Train surrogate locally, craft attacks offline, transfer to production API. No queries needed.}
    \end{mlsyscard}
  \end{column}
  \begin{column}{0.48\textwidth}
    \begin{mlsyscard}{routingstroke}
    {\scriptsize \textbf{Physical Patches}: Adversarial patches on clothing, stickers on stop signs. No digital access required.}
    \end{mlsyscard}
  \end{column}
\end{columns}

\end{frame}

% --- ACTIVE LEARNING 2: Exercise ---
\begin{frame}{Your Turn: Diagnose the Threat}
\note{
% -- LINK: What prior concept connects to this slide
Students learned drift types and adversarial attacks. This exercise tests whether
they can diagnose a real-world scenario combining both concepts.

% -- NARRATE: What to SAY while showing this slide
[3 min] Read the scenario. Give 90 seconds. After pause reveal: ``This is concept
drift --- P(Y|X) changed because fraud patterns evolved. Input features are
unchanged, so KS test shows no signal. You need ground-truth comparison.''

% -- ENGAGE: Specific question, prediction, or task for THIS slide
Three-part question forces precise reasoning: identify the shift type, explain
why monitoring missed it, propose the correct response. [Expected: concept drift,
invisible to input monitoring, retrain every 90--120 days.]

% -- FLEX: [CORE] or [OPTIONAL] + contingency
[CORE] Second required active learning moment.
IF SHORT: Reduce to 60 seconds; skip pair comparison.
}

\footnotesize
\begin{columns}[T]
  \begin{column}{0.55\textwidth}
    {\small\bfseries Diagnose the Threat}

    \vspace{0.1cm}
    A fraud detection model deployed 6 months ago:
    \begin{itemize}\setlength\itemsep{0pt}
      \item Input features: \textbf{unchanged}
      \item Model confidence: \textbf{stable at 95\%+}
      \item Fraud precision: \textbf{85\% $\to$ 68\%}
    \end{itemize}

    \vspace{0.1cm}
    \textbf{1.} What type of shift is this?\\
    \textbf{2.} Why did input monitoring miss it?\\
    \textbf{3.} What is the correct response?

    {\scriptsize\textcolor{midgray}{(90 seconds --- compare with a neighbor)}}
  \end{column}
  \begin{column}{0.42\textwidth}
    \pause
    \begin{mlsyscard}{datastroke}
    \textbf{Solution:}\\[0.05cm]
    {\scriptsize
    \textbf{1.} Concept drift: $P(Y|X)$ changed (fraud patterns evolved)\\[0.05cm]
    \textbf{2.} $P(X)$ unchanged --- input monitoring only detects covariate/label shift\\[0.05cm]
    \textbf{3.} Ground-truth feedback loop + retrain every 90--120 days\\[0.05cm]
    \textcolor{errorstroke}{\textbf{Concept drift is invisible to feature monitoring!}}
    }
    \end{mlsyscard}
  \end{column}
\end{columns}

\end{frame}

% =============================================================================
\section{Defenses}
% =============================================================================

\begin{frame}{The Robustness Tax}
\note{
% -- LINK: What prior concept connects to this slide
We have cataloged threats. Now the defense section opens with the uncomfortable
truth: every defense has a cost.

% -- NARRATE: What to SAY while showing this slide
[3 min] Point to the bar chart: ``Standard ResNet-50: 76\% accuracy. Adversarial
training: 50\% --- a 26\% tax on clean accuracy. Certified smoothing: requires
100K noise samples per inference.'' Key punchline: ``Robustness cannot be turned
on for free. Every defense trades accuracy, compute, or both.''

% -- WARN: What students will get wrong on THIS topic
Common error: students assume robustness is additive --- just add a defense layer.
Correct framing: adversarial training fundamentally changes what the model learns,
sacrificing non-robust features that were predictive on clean data.

% -- FLEX: [CORE] or [OPTIONAL] + contingency
[CORE] The robustness tax frames all subsequent defense slides.
IF SHORT: State the 26\% accuracy tax and the 100K inference cost.
}

% --- Layout: FULL-WIDTH diagram ---
\centering
\safeimg[width=0.92\textwidth,height=4.8cm,keepaspectratio]{robustness-tax.pdf}

\vspace{0.1cm}
\mlsysalert{Fundamental:}{Robustness trades accuracy, compute, or both. Every defense has a cost.}

\end{frame}

\begin{frame}{Adversarial Training}
\note{
% -- LINK: What prior concept connects to this slide
The robustness tax showed the cost. This slide explains the two primary defense
mechanisms and why each costs what it does.

% -- NARRATE: What to SAY while showing this slide
[2 min] Point to the left column: ``Adversarial training augments each batch with
PGD-generated perturbations. Cost: 8--10x compute per epoch. Clean accuracy drops
from 76\% to 50\%.'' Point to the right column: ``Certified smoothing provides
mathematical guarantees via randomized noise. But 100K samples per inference
restricts it to offline auditing.''
ANALOGY: ``Adversarial training is like sparring practice --- expensive but builds
real resilience. Certified smoothing is like armor testing --- rigorous but slow.''

% -- WARN: What students will get wrong on THIS topic
Common error: students think adversarial training protects against all attacks.
Correct framing: it is threat-model-specific --- L-infinity training provides
zero protection against L-2 or geometric attacks.

% -- FLEX: [CORE] or [OPTIONAL] + contingency
[CORE] These are the two primary defense paradigms.
IF SHORT: Cover adversarial training only; mention certified smoothing as offline.
}

\small
\begin{columns}[T]
  \begin{column}{0.48\textwidth}
    \textbf{Empirical Defense:}

    \vspace{0.1cm}
    Adversarial training augments each batch with PGD-generated perturbations:

    \vspace{0.1cm}
    {\footnotesize
    \begin{itemize}\setlength\itemsep{0pt}
      \item Generate adversarial $x_{\text{adv}}$ for each $x$
      \item Train on both clean and adversarial
      \item Cost: \textbf{8--10$\times$} compute per epoch
      \item Clean accuracy: 76\% $\to$ \textbf{50\%}
    \end{itemize}
    }

    \vspace{0.1cm}
    \mlsysconcept{Why it works:}{Model learns to ignore brittle, non-robust features.}
  \end{column}
  \begin{column}{0.48\textwidth}
    \textbf{Certified Defense:}

    \vspace{0.1cm}
    Randomized smoothing provides \emph{mathematical} guarantees:

    \vspace{0.1cm}
    {\footnotesize
    $g(x) = \arg\max_c \mathbb{P}(f(x+\delta) = c)$\\
    where $\delta \sim \mathcal{N}(0, \sigma^2 I)$
    }

    \vspace{0.1cm}
    {\footnotesize
    \begin{itemize}\setlength\itemsep{0pt}
      \item Certified radius: $R = \sigma \Phi^{-1}(p_A)$
      \item 49\% certified robustness
      \item Cost: \textbf{100K samples/inference}
      \item Restricted to offline auditing
    \end{itemize}
    }
  \end{column}
\end{columns}

\end{frame}

\begin{frame}{Data Poisoning: Attacking the Supply Chain}
\note{
% -- LINK: What prior concept connects to this slide
Adversarial training defends inference. Data poisoning attacks training ---
a different phase requiring different defenses (also covered in Ch 13).

% -- NARRATE: What to SAY while showing this slide
[3 min] Point to the diagram: ``Clean-label attacks are the most insidious:
labels are correct, but images contain invisible triggers. 0.1\% poisoned
data achieves 95\% backdoor success.'' Key point: ``No pipeline access needed
--- just contribute to a public dataset like ImageNet or LAION.''

% -- WARN: What students will get wrong on THIS topic
Common error: students think poisoning requires admin access to training data.
Correct framing: clean-label attacks embed triggers in publicly contributed
data --- anyone can poison a crowdsourced dataset.

% -- FLEX: [CORE] or [OPTIONAL] + contingency
[CORE] Supply chain attacks are the highest-leverage training-time threat.
IF SHORT: State the 0.1\% number and the clean-label insight.
}

% --- Layout: FULL-WIDTH diagram ---
\centering
\safeimg[width=0.92\textwidth,height=4.5cm,keepaspectratio]{data-poisoning.pdf}

\vspace{0.1cm}
\mlsysalert{Supply chain attack:}{0.1\% poisoned training data embeds a 95\%-effective backdoor.}

\end{frame}

\begin{frame}{Multi-Layered Defense}
\note{
% -- LINK: What prior concept connects to this slide
Individual defenses each have blind spots. This slide assembles them into a
layered defense where each ring catches what the others miss.

% -- NARRATE: What to SAY while showing this slide
[3 min] Walk through the concentric rings: ``Outer ring: feature squeezing ---
eliminates 70--90\% of adversarial examples at negligible cost. Next: adversarial
training provides worst-case guarantees. Drift monitoring gives 3--6 weeks of
early warning. Center: uncertainty quantification --- the model should be able
to say `I do not know.'\,''

% -- ENGAGE: Specific question, prediction, or task for THIS slide
Ask: ``Which layer provides the best cost-to-value ratio?'' Give 10 seconds.
[Expected: feature squeezing --- 70--90\% effectiveness at negligible cost.]

% -- FLEX: [CORE] or [OPTIONAL] + contingency
[CORE] Defense-in-depth is the chapter's design pattern.
IF SHORT: Name the layers and highlight feature squeezing as highest leverage.
}

% --- Layout: FULL-WIDTH diagram ---
\centering
\safeimg[width=0.88\textwidth,height=4.8cm,keepaspectratio]{defense-layers.pdf}

\vspace{0.1cm}
\mlsysinsight{Defense-in-depth:}{Each layer catches failures the others miss. No single technique is sufficient.}

\end{frame}

\begin{frame}{Detection Techniques}
\note{
% -- LINK: What prior concept connects to this slide
The layered defense diagram showed the strategy. This slide provides the specific
detection techniques that populate those layers.

% -- NARRATE: What to SAY while showing this slide
[2 min] Walk through the table: ``Statistical methods compare input to training
distribution. Feature squeezing reduces precision from 256 to 16 colors ---
70--90\% of adversarial examples filtered at negligible latency. MC dropout runs
multiple forward passes for uncertainty estimates at 10--100x inference cost.''
Point to the green card: ``Feature squeezing is the highest value-to-cost ratio.''

% -- FLEX: [OPTIONAL] or [CORE] + contingency
[OPTIONAL] Extends the defense layers slide with implementation detail.
IF SHORT: Point to feature squeezing in the green card and move on.
}

\scriptsize
\renewcommand{\arraystretch}{1.0}
\begin{tabular}{@{}lp{2.8cm}p{2cm}p{2.8cm}@{}}
  \toprule
  \textbf{Approach} & \textbf{Mechanism} & \textbf{Effectiveness} & \textbf{Cost} \\
  \midrule
  \textbf{Statistical}    & Compare input to training dist.\ & KS, AD tests & Low (feature-level) \\
  \textbf{Feature Squeeze} & Reduce precision (256$\to$16) & 70--90\% filtered & Negligible latency \\
  \textbf{MC Dropout}      & Multiple forward passes & Uncertainty est.\ & 10--100$\times$ inference \\
  \textbf{Ensembles}       & Multiple models, variance & Robust UQ & 3--5$\times$ compute \\
  \bottomrule
\end{tabular}

\vspace{0.1cm}
\begin{mlsyscard}{datastroke}
{\scriptsize \textbf{Feature Squeezing} (highest value / lowest cost): Compare predictions on original vs.\ squeezed inputs. Large divergence flags adversarial. 95\%+ clean accuracy at negligible latency.}
\end{mlsyscard}

\end{frame}

% --- ACTIVE LEARNING 3: Discussion ---
\begin{frame}{Discussion: Choose Your Defense}
\note{
% -- LINK: What prior concept connects to this slide
Students just learned five defense techniques with different cost profiles. This
discussion forces them to choose under real constraints.

% -- NARRATE: What to SAY while showing this slide
[3 min] Read the scenario: ``Medical imaging classifier. Budget: 2x training
compute, 1.5x inference latency.'' Point to the five options. Give 90 seconds
for pair discussion. Cold-call 2--3 pairs.

% -- ENGAGE: Specific question, prediction, or task for THIS slide
Students must choose AND justify. [No single right answer --- the point is that
defense choice depends on threat model and compute budget. Adversarial training
fits the training budget but exceeds inference; feature squeezing fits inference
but provides weaker guarantees.]

% -- FLEX: [CORE] or [OPTIONAL] + contingency
[CORE] Third required active learning moment.
IF SHORT: Reduce to 60 seconds; cold-call 2 pairs.
}

\centering
\vspace{0.6cm}
{\large\bfseries Turn and Talk \textcolor{midgray}{(90 seconds)}}

\vspace{0.5cm}
{\large You are deploying a medical imaging classifier.\\
Budget: 2$\times$ training compute, 1.5$\times$ inference latency.\\[0.3cm]
\alert{Which defense strategy do you choose and why?}}

\vspace{0.5cm}
\begin{columns}[c]
  \begin{column}{0.20\textwidth}\centering\small Adversarial\\ Training\end{column}
  \begin{column}{0.20\textwidth}\centering\small Certified\\ Smoothing\end{column}
  \begin{column}{0.20\textwidth}\centering\small Feature\\ Squeezing\end{column}
  \begin{column}{0.20\textwidth}\centering\small Ensemble +\\ MC Dropout\end{column}
  \begin{column}{0.20\textwidth}\centering\small Drift\\ Monitoring\end{column}
\end{columns}

\end{frame}

% =============================================================================
\section{Trade-offs}
% =============================================================================

\begin{frame}{The Robustness--Efficiency Tension}
\note{
% -- LINK: What prior concept connects to this slide
Defenses have costs. This slide quantifies the aggregate overhead and connects
robustness to sustainability (Ch 15).

% -- NARRATE: What to SAY while showing this slide
[2 min] Walk through the overhead table: ``Error correction: 12--25\% bandwidth.
Redundant processing: 2--3x energy. Monitoring: 5--15\% compute. Adversarial
training: 8--10x per epoch.'' Point to the crimson card: ``Robustness is a
resource allocation problem. A system that survives every failure yet consumes
a city's worth of power is not viable.''

% -- ENGAGE: Specific question, prediction, or task for THIS slide
Ask: ``Where does robustness investment provide the greatest value --- cloud,
edge, or TinyML?'' [Expected: depends on criticality of the application.]

% -- FLEX: [CORE] or [OPTIONAL] + contingency
[CORE] Connects robustness costs to the sustainability chapter.
IF SHORT: State the key overhead numbers from the table.
}

\footnotesize
\begin{columns}[T]
  \begin{column}{0.48\textwidth}
    \textbf{Robustness costs are real:}

    \vspace{0.05cm}
    \renewcommand{\arraystretch}{1.0}
    {\scriptsize
    \begin{tabular}{@{}lr@{}}
      \toprule
      \textbf{Mechanism} & \textbf{Overhead} \\
      \midrule
      Error correction    & 12--25\% BW \\
      Redundant processing & 2--3$\times$ energy \\
      Monitoring           & 5--15\% compute \\
      Adv.\ training       & 8--10$\times$/epoch \\
      Certified inference  & 100K$\times$ latency \\
      \bottomrule
    \end{tabular}
    }

    \vspace{0.05cm}
    {\scriptsize\mlsysalert{Decision:}{Where does robustness investment provide the greatest value?}}
  \end{column}
  \begin{column}{0.48\textwidth}
    \begin{mlsyscard}{crimson}
    \textbf{The Core Trade-off}\\[0.05cm]
    {\scriptsize Robustness is a resource allocation problem. A system that survives every failure yet consumes a city's worth of power is not viable.}
    \end{mlsyscard}

    \vspace{0.1cm}
    \begin{mlsyscard}{routingstroke}
    \textbf{Deploy-Tier Matching}\\[0.05cm]
    {\scriptsize Cloud: full redundancy + monitoring.\\
    Edge: targeted hardening of critical paths.\\
    TinyML: weaker guarantees on non-critical paths.}
    \end{mlsyscard}
  \end{column}
\end{columns}

\end{frame}

\begin{frame}{Compound Threats}
\note{
% -- LINK: What prior concept connects to this slide
Individual threats were covered separately. The final insight: in production,
they arrive together and interact.

% -- NARRATE: What to SAY while showing this slide
[2 min] Walk through the bullet list: ``A software fault creates a new adversarial
vulnerability. A distribution shift masks a subtle poisoning attack. Quantization
to INT8 reduces robustness margins. A pipeline bug mimics an adversarial attack
in monitoring.'' Key statistic: ``60\%+ of ML incidents stem from pipeline and
data issues, not models.''

% -- WARN: What students will get wrong on THIS topic
Common error: students analyze threats in isolation. Correct framing: compound
threats are multiplicative, not additive --- a shift that also triggers a
software bug is worse than either alone.

% -- FLEX: [CORE] or [OPTIONAL] + contingency
[CORE] Compound threats justify the multi-layered defense approach.
IF SHORT: State the 60\% pipeline statistic and the compounding principle.
}

\footnotesize
\begin{columns}[T]
  \begin{column}{0.55\textwidth}
    \textbf{Failures compound, not just add:}

    \vspace{0.05cm}
    {\scriptsize
    \begin{itemize}\setlength\itemsep{0pt}
      \item \textcolor{computestroke}{\textbf{Software fault}} creates new adversarial vulnerability
      \item \textcolor{routingstroke}{\textbf{Distribution shift}} masks subtle data poisoning
      \item \textcolor{errorstroke}{\textbf{Quantization}} (INT8) reduces robustness margin
      \item \textcolor{datastroke}{\textbf{Pipeline bug}} mimics adversarial attack in monitoring
    \end{itemize}
    }

    \vspace{0.05cm}
    {\scriptsize\mlsysconcept{Key:}{60\%+ of ML incidents stem from pipeline/data issues, not models.}}
  \end{column}
  \begin{column}{0.42\textwidth}
    \begin{mlsyscard}{errorstroke}
    {\scriptsize \textbf{Defense Implication}: No single technique addresses compound threats. Combine:
    \begin{enumerate}\setlength\itemsep{0pt}
      \item Input sanitization
      \item Robust training objectives
      \item Output verification
      \item Continuous monitoring
      \item Uncertainty quantification
    \end{enumerate}
    }
    \end{mlsyscard}
  \end{column}
\end{columns}

\end{frame}

% =============================================================================
\section{Wrap-Up}
% =============================================================================

\begin{frame}{Fallacies}
\note{
% -- NARRATE: What to SAY while showing this slide
[2 min] Read each fallacy and its rebuttal. Emphasize quantitative evidence:
physical patches on stop signs, 20--40\% OOD accuracy drops, 0.1\% poisoning
from public datasets, L-infinity training offering zero L-2 protection.

% -- FLEX: [CORE] or [OPTIONAL] + contingency
[CORE] Fallacies correct the most dangerous misconceptions.
IF SHORT: Read fallacies 1 and 4 (academic curiosity and adversarial training).
}

\footnotesize
\textbf{Fallacy:} \textit{Adversarial examples are an academic curiosity.}\\
{\scriptsize Physical-world patches on stop signs reliably fool production vision systems. Defense requires 10--20\% additional compute.}

\vspace{0.05cm}
\textbf{Fallacy:} \textit{If the model passes the test set, it is robust.}\\
{\scriptsize I.I.D.\ test sets cannot measure OOD resilience. Drift causes 20--40\% accuracy drops on out-of-distribution inputs.}

\vspace{0.05cm}
\textbf{Fallacy:} \textit{Data poisoning requires pipeline access.}\\
{\scriptsize Clean-label attacks inject malicious samples into \emph{public} datasets. 0.1\% poisoned data $\to$ 95\% backdoor success.}

\vspace{0.05cm}
\textbf{Fallacy:} \textit{Adversarial training solves the problem.}\\
{\scriptsize Robustness is threat-model-specific. $L_\infty$-trained models offer zero protection against $L_2$ or geometric attacks.}

\end{frame}

\begin{frame}{Pitfalls}
\note{
% -- NARRATE: What to SAY while showing this slide
[2 min] Read each pitfall. The first is critical: deploying without drift
monitoring means silent degradation for weeks. The second: average accuracy
hides 100\% vulnerability on critical edge cases. The third: 60\%+ of ML
incidents are pipeline issues, not model failures.

% -- FLEX: [CORE] or [OPTIONAL] + contingency
[CORE] Pitfalls are operational guidance students will use immediately.
IF SHORT: Cover the first two pitfalls only.
}

\footnotesize
\textbf{Pitfall:} \textit{Deploying without distribution shift monitoring.}\\
Models degrade silently while maintaining high confidence scores. PSI monitoring detects shifts \textbf{3--6 weeks} before accuracy falls below SLA thresholds.

\vspace{0.15cm}
\textbf{Pitfall:} \textit{Using average accuracy as the sole robustness metric.}\\
95\% average accuracy can coexist with 100\% vulnerability on critical edge cases. Evaluate \textbf{worst-case accuracy} under a specific perturbation budget ($\eps = 8/255$).

\vspace{0.15cm}
\textbf{Pitfall:} \textit{Ignoring software faults because the model is correct.}\\
Over \textbf{60\%} of ML incidents stem from pipeline and data issues, not models. A model hardened against adversarial attacks remains fragile if a preprocessing bug silently corrupts its inputs.

\end{frame}

% --- RETRIEVAL PRACTICE ---

% --- MUDDIEST POINT ---
\begin{frame}{Muddiest Point}
\note{
% -- NARRATE: What to SAY while showing this slide
[2 min] ``Write down the single concept from today you found most confusing.
Anonymous, one sentence.'' Collect responses. Scan for patterns --- likely
candidates: robustness tax math, concept drift detection, or compound threats.

% -- FLEX: [CORE] + contingency
[CORE] Closes the feedback loop. Do not skip.
}

\centering
\vspace{1.0cm}
{\Large\bfseries What was the \alert{muddiest point} today?}

\vspace{0.8cm}
{\normalsize Write down the concept you found \textbf{most confusing}.}

\vspace{0.5cm}
{\small\textcolor{midgray}{Anonymous. One sentence. Submit before you leave.}}

\end{frame}

\begin{frame}{What Were the Key Ideas?}
\note{
% -- NARRATE: What to SAY while showing this slide
[2 min] ``Close your notes. Write down the 4 most important concepts from today.
90 seconds, no peeking.'' Walk around the room. Do NOT show next slide yet.

% -- FLEX: [CORE] + contingency
[CORE] Retrieval practice is the highest-impact learning intervention.
}

\centering
\vspace{1.5cm}
{\Large\bfseries Close your notes.}

\vspace{0.8cm}
{\large Write down the \textbf{4 most important concepts} from today.}

\vspace{0.8cm}
{\normalsize\textcolor{midgray}{90 seconds --- no peeking.}}

\end{frame}

\begin{frame}{Key Takeaways}
\note{
% -- NARRATE: What to SAY while showing this slide
[2 min] Reveal after retrieval. Walk through each bullet, anchoring on numbers:
silent 15--25\% degradation, 8--10x adversarial training cost, 26\% clean accuracy
tax, 3--6 week PSI early warning, 60\%+ pipeline incidents.

% -- FLEX: [CORE] + contingency
[CORE] Synthesis slide. IF SHORT: Read bullets 1, 3, and 5.
}

\footnotesize
\begin{itemize}\setlength\itemsep{1pt}
  \item \textbf{Silent failure is the real threat}: ML systems degrade without crashes, alerts, or exceptions. Confidence stays high while accuracy drops 15--25\% over months.
  \item \textbf{Robustness $\neq$ accuracy}: A model with 95\% test accuracy can be 100\% vulnerable to $\eps = 8/255$ perturbations. Standard test sets measure i.i.d.\ performance, not robustness.
  \item \textbf{The robustness tax is real}: Adversarial training costs 8--10$\times$ compute and 26\% clean accuracy. Certified smoothing: 100K$\times$ inference latency.
  \item \textbf{Three shift types}: Covariate (P(X) changes), label (P(Y) changes), concept (P(Y$|$X) changes). Only concept drift requires ground-truth feedback.
  \item \textbf{Multi-layered defense}: Input sanitization, adversarial training, drift monitoring, data validation, and uncertainty quantification. Each catches what the others miss.
  \item \textbf{Compound threats}: Software faults, distribution shift, and adversarial attacks interact and amplify each other. 60\%+ of incidents are pipeline/data issues.
\end{itemize}

\end{frame}

\begin{frame}{References}
\note{
% -- NARRATE: What to SAY while showing this slide
[1 min] Highlight Madry+18 as the adversarial training foundation and Dixit+21
for silent data corruption at scale. Students can photograph for later reading.

% -- FLEX: [OPTIONAL] + contingency
[OPTIONAL] Do not read every entry aloud.
}

\small
\mlsysref{Goodfellow+15}{Goodfellow et al. ``Explaining and Harnessing Adversarial Examples.'' ICLR 2015.}
\mlsysref{Madry+18}{Madry et al. ``Towards Deep Learning Models Resistant to Adversarial Attacks.'' ICLR 2018.}
\mlsysref{Carlini+17}{Carlini \& Wagner. ``Towards Evaluating the Robustness of Neural Networks.'' IEEE S\&P 2017.}
\mlsysref{Cohen+19}{Cohen et al. ``Certified Adversarial Robustness via Randomized Smoothing.'' ICML 2019.}
\mlsysref{Dixit+21}{Dixit et al. ``Silent Data Corruptions at Scale.'' arXiv 2021.}
\mlsysref{NTSB17}{NTSB. ``Tesla Autopilot Fatal Crash Investigation.'' 2017.}

\end{frame}

\begin{frame}{Next Lecture: Sustainable AI}
\note{
% -- NARRATE: What to SAY while showing this slide
[1 min] ``We made the system robust --- but at what cost? Adversarial training,
redundant processing, continuous monitoring all consume energy.'' Point to the
three cards. ``If the fleet survives every failure but consumes a city's worth
of power, that is not viable. Next lecture: sustainability as a hard constraint.''

% -- FLEX: [CORE] + contingency
[CORE] Forward hook motivates attendance at the next lecture.
}

\footnotesize
\begin{columns}[c]
  \begin{column}{0.30\textwidth}
    \centering
    \begin{mlsyscard}{computestroke}
    {\large\bfseries Carbon}\\[0.1cm]
    {\footnotesize Lifecycle carbon accounting for ML training and inference}
    \end{mlsyscard}
  \end{column}
  \begin{column}{0.30\textwidth}
    \centering
    \begin{mlsyscard}{datastroke}
    {\large\bfseries Efficiency}\\[0.1cm]
    {\footnotesize Algorithmic efficiency and hardware utilization optimization}
    \end{mlsyscard}
  \end{column}
  \begin{column}{0.30\textwidth}
    \centering
    \begin{mlsyscard}{routingstroke}
    {\large\bfseries Scheduling}\\[0.1cm]
    {\footnotesize Carbon-aware job scheduling and renewable energy alignment}
    \end{mlsyscard}
  \end{column}
\end{columns}

\vspace{0.2cm}
\centering
{\normalsize A system that survives every failure yet consumes a city's worth of power}\\[0.1cm]
{\small\textbf{is not a viable long-term solution.}}

\end{frame}



\appendix

\begin{frame}{Backup: Extended Reference}
\note{
% -- NARRATE: What to SAY while showing this slide
Backup. Reference the textbook for FGSM derivation details, certified smoothing
proofs, and additional SDC case studies from Meta and Google.

% -- FLEX: [OPTIONAL] + contingency
[OPTIONAL] Only show if students ask for deeper material.
}

\footnotesize
This slide provides extended reference material for students who want to go deeper.
Refer to the textbook chapter for complete derivations and additional worked examples.

\vspace{0.3cm}
\begin{mlsyscard}{crimson}
\textbf{Key equations and frameworks from this chapter} are collected in the
textbook's summary tables. Use them as a quick reference during problem sets.
\end{mlsyscard}

\end{frame}

\begin{frame}{Backup: Further Reading}
\note{
% -- NARRATE: What to SAY while showing this slide
Backup. Point students to the textbook's extended adversarial attack taxonomy,
the drift monitoring implementation guide, and the compound threats analysis.

% -- FLEX: [OPTIONAL] + contingency
[OPTIONAL] Only show if time permits after Q\&A.
}

\footnotesize
\textbf{For deeper exploration:}

\vspace{0.15cm}
\begin{itemize}\setlength\itemsep{2pt}
  \item The textbook chapter contains 2--3 additional worked examples
  \item Problem sets apply these concepts to real-world scenarios
  \item The course website links to open-source implementations
  \item Office hours: bring your toughest questions
\end{itemize}

\end{frame}


\end{document}
